\documentclass[12pt,a4paper]{article}
\usepackage[utf8]{inputenc}
\usepackage[french]{babel}
\usepackage{geometry}
\geometry{margin=2cm}
\usepackage{amsmath,amssymb,amsthm}
\usepackage{graphicx}
\usepackage{booktabs}
\usepackage{array}
\usepackage{xcolor}
\usepackage{hyperref}
\usepackage{float}
\usepackage{parskip}
\usepackage[]{multicol}
\title{\textbf{RAPPORT DE SIMULATION STRUCTURELLE - ANSYS}}
\author{Friedly Woli \\ Étudiant L3 Mécanique Énergétique \\ Université Toulouse III Paul Sabatier}
\date{06 Mars 2026}

\begin{document}
	
	\maketitle
	\tableofcontents
	\newpage
	
	\section{CONTEXTE DE LA SIMULATION}
	
			\subsection{Projet}
			\begin{itemize}
				\item \textbf{Contexte projet :} \dotfill % À compléter
				\item \textbf{Organisation :} \dotfill % Tolosat SUPAERO / TLS'E Racing
				\item \textbf{Rôle personnel :} \dotfill % Division Structure / Powertrain 
			\end{itemize}

			\subsection{Objectifs de la simulation}

			
				
				
				\begin{multicols}{2}
					\textbf{ Objectifs } \\
					Verifier la resistance et controler la deformation, Identifier la deformation 
					
					\columnbreak
					
					\textbf{Criteres de validation }\\
					$\sigma_{VM} < \sigma_{limite}$ \\
					$u_{max} < u_{admise}$ \\
					Localisation précise des points chauds \\
					Variation $< 5\%$ entre maillages \\
				\end{multicols}
				
			
	\newpage
	
	
	
	
	\section{MODÉLISATION}
	
					\subsection{Caractéristiques géométriques}
					\begin{table}[H]
						\centering
						\begin{tabular}{lll}
							\toprule
							\textbf{Paramètre} & \textbf{Valeur} & \textbf{Unité} \\
							\midrule
							Dimensions globales & \dotfill & mm \\
							Masse de la pièce & \dotfill & kg \\
							Symétries exploitées & \dotfill & \text{(Oui/Non)} \\
							Simplifications géométriques & \dotfill & \text{(trous, dépouilles...)} \\
							\bottomrule
						\end{tabular}
						\caption{Caractéristiques géométriques}
					\end{table}
			
					\subsection{Matériau}
					\begin{table}[H]
						\centering
						\begin{tabular}{p{4cm}p{3cm}p{3cm}p{5cm}}
							\toprule
							\textbf{Propriété} & \textbf{Valeur} & \textbf{Unité} & \textbf{Source} \\
							\midrule
							Module d'Young & \dotfill & GPa & \dotfill \\
							Coefficient Poisson & \dotfill & - & \dotfill \\
							Limite élastique & \dotfill & MPa & \dotfill \\
							Densité & \dotfill & kg/m$^3$ & \dotfill \\
							Conductivité thermique & \dotfill & W/m.K & \dotfill \\
							\bottomrule
						\end{tabular}
						\caption{Caractéristiques du matériau}
					\end{table}
			
					\subsection{Conditions aux limites}
					\begin{table}[H]
						\centering
						\begin{tabular}{p{3cm}p{4cm}p{5cm}p{4cm}}
							\toprule
							\textbf{Type} & \textbf{Emplacement} & \textbf{Valeur} & \textbf{Justification} \\
							\midrule
							Encadrement & \dotfill & \dotfill & \dotfill \\
							Liaison élastique & \dotfill & \dotfill & \dotfill \\
							Charge ponctuelle & \dotfill & \dotfill & \dotfill \\
							Pression & \dotfill & \dotfill & \dotfill \\
							Moment & \dotfill & \dotfill & \dotfill \\
							\bottomrule
						\end{tabular}
						\caption{Conditions aux limites}
					\end{table}
			
			\subsection{ÉTUDE DE CONVERGENCE DU MAILLAGE}
			
					\begin{table}[H]
						\centering
						\begin{tabular}{ccccc}
							\toprule
							\textbf{Maillage} & \textbf{Nbre éléments} & \textbf{$u_{max}$ (mm)} & \textbf{$\sigma_{VM,max}$ (MPa)} & \textbf{Variation (\%)} \\
							\midrule
							Grossier & \dotfill & 2.50 & 280 & - \\
							Intermédiaire & \dotfill & 2.30 & 265 & 8.0 \\
							Fin & \dotfill & 2.28 & 258 & 1.3 \\
							Très fin & \dotfill & 2.27 & 257 & 0.4 \\
							\bottomrule
						\end{tabular}
						\caption{Étude de convergence}
					\end{table}
					
					\textbf{Conclusion convergence :} Maillage \textit{Fin} retenu (variation $< 2\%$).
			
			
			
			
			
			
\newpage






			\section{RÉSULTATS}
			
					\subsection{Déformation totale}
					\begin{table}[H]
						\centering
						\begin{tabular}{lcc}
							\toprule
							\textbf{Paramètre} & \textbf{Valeur} & \textbf{Limite admissible} \\
							\midrule
							Déformation max & \dotfill mm & \dotfill mm \\
							Déformation moyenne & \dotfill mm & - \\
							Zone critique & \dotfill & - \\
							Facteur de sécurité & \dotfill & $>1.5$ \\
							\bottomrule
						\end{tabular}
						\caption{Résultats déformation}
					\end{table}
					
					\subsection{Contraintes Von Mises}
					\begin{table}[H]
						\centering
						\begin{tabular}{lcc}
							\toprule
							\textbf{Paramètre} & \textbf{Valeur} & \textbf{Limite admissible} \\
							\midrule
							$\sigma_{VM,max}$ & \dotfill MPa & \dotfill MPa \\
							$\sigma_{VM,moyenne}$ & \dotfill MPa & - \\
							Zone critique & \dotfill & - \\
							Facteur de sécurité & \dotfill & $>1.5$ \\
							\bottomrule
						\end{tabular}
						\caption{Résultats contraintes Von Mises}
					\end{table}
					
					\begin{figure}[H]
						\centering
					%	\includegraphics[width=0.45\textwidth]{deformation-total1.jpg}
					%	\includegraphics[width=0.45\textwidth]{equivalent-stress.jpg}
						\caption{Déformation totale (gauche) et contraintes Von Mises (droite)}
					\end{figure}
			
			
			
			
			
			
	\newpage
	
	
	
	
	
	
	
	\section{VALIDATION DES OBJECTIFS}
	
					\begin{table}[H]
						\centering
						\begin{tabular}{p{3cm}|p{4cm}|p{4cm}|p{3cm}}
							\toprule
							\textbf{Objectif} & \textbf{Résultat obtenu} & \textbf{Critère respecté ?} & \textbf{Justification} \\
							\midrule
							Résistance & $\sigma_{VM,max} = $ \dotfill MPa & \dotfill & \dotfill \\
							Déformation & $u_{max} = $ \dotfill mm & \dotfill & \dotfill \\
							Convergence & Variation $< 2\%$ & \textbf{OUI} & Tableau convergence \\
							Zones critiques & Identifiées aux coins & \textbf{OUI} & Images jointes \\
							\bottomrule
						\end{tabular}
						\caption{Validation des objectifs}
					\end{table}
			
				\subsection{ANALYSE DES ZONES CRITIQUES}
				
						\begin{enumerate}
							\item \textbf{Coin interne 1 :} $\sigma_{VM} = $ \dotfill MPa $\to$ \dotfill
							\item \textbf{Coin interne 2 :} $\sigma_{VM} = $ \dotfill MPa $\to$ \dotfill
							\item \textbf{Coin interne 3 :} $\sigma_{VM} = $ \dotfill MPa $\to$ \dotfill
							\item \textbf{Bords externes :} Faibles contraintes $\to$ \textbf{Possibilité d'allègement}
						\end{enumerate}
				
				\subsection{RECOMMANDATIONS DE CONCEPTION}
				
						\begin{table}[H]
							\centering
							\begin{tabular}{p{6cm}p{6cm}}
								\toprule
								\textbf{Action corrective} & \textbf{Impact attendu} \\
								\midrule
								Renforcer coins internes (raidisseurs) & $\sigma_{VM} ↓ 20-30\%$ \\
								Augmenter épaisseur locale & $u_{max} ↓ 15\%$ \\
								Alléger zones bleues & Masse ↓ 10\% sans perte rigidité \\
								Modifier appui vert & Répartition homogène \\
								Re-simulation après modifs & Validation finale \\
								\bottomrule
							\end{tabular}
							\caption{Recommandations d'amélioration}
						\end{table}
			
			
			
			
			
			
	\newpage
	
	
	
	
	
	
	\section{CONCLUSION}
	
			\begin{itemize}
				\item \textbf{Simulation validée} : convergence atteinte, objectifs remplis
				\item \textbf{Zones critiques identifiées} : 3 coins internes à traiter
				\item \textbf{Facteurs de sécurité} : \dotfill (acceptable / à améliorer)
				\item \textbf{Prochaines étapes :} \dotfill (modification géométrie / prototype / essais)
			\end{itemize}
			
			\vspace{1cm}
			\begin{center}
				\textbf{Rapport validé le : \dotfill}
			\end{center}
	
\end{document}
